\section{Derivation: Viewing SFT as a Special Case of On-Policy RL}
\label{sec:proof}

In this appendix, we provide a detailed derivation showing that supervised fine-tuning (SFT) can be interpreted as a special case of reinforcement learning (RL) with a sparse reward function. Specifically, we show that the gradient of the SFT objective can be rewritten as an on-policy expectation under the current policy via importance sampling.

\subsection{SFT Objective and Gradient}

We consider a dataset of expert demonstrations
\(
\mathcal{D} = \{(x, y^*)\}
\),
where \(x\) denotes the input and \(y^*\) is the expert-provided output. The standard SFT objective is defined as the negative log-likelihood:
\begin{equation}
\mathcal{L}_{\mathrm{SFT}}(\theta)
=
-\mathbb{E}_{(x,y^*) \sim \mathcal{D}}
\left[
\log \pi_\theta(y^* \mid x)
\right].
\end{equation}

Taking the gradient with respect to the model parameters \(\theta\), we obtain
\begin{equation}
\nabla_\theta \mathcal{L}_{\mathrm{SFT}}(\theta)
=
-
\mathbb{E}_{(x,y^*) \sim \mathcal{D}}
\left[
\nabla_\theta \log \pi_\theta(y^* \mid x)
\right].
\label{eq:sft_grad_expert}
\end{equation}

This expectation is taken over the expert data distribution rather than samples generated by the current policy.

\subsection{Importance Sampling Reformulation}

We factorize the expert data distribution as
\begin{equation}
P(x, y^*) = P(x)\, P_{\mathrm{expert}}(y^* \mid x),
\end{equation}
and define the joint distribution induced by the current policy as
\begin{equation}
Q(x, y) = P(x)\, \pi_\theta(y \mid x).
\end{equation}

Since both distributions share the same marginal \(P(x)\), we can apply importance sampling to rewrite the expectation in Eq.~\eqref{eq:sft_grad_expert} under \(Q(x,y)\):
\begin{equation}
\begin{aligned}
&\nabla_\theta \mathcal{L}_{\mathrm{SFT}}(\theta)\\
&=
-
\mathbb{E}_{(x,y) \sim Q}
\left[
\nabla_\theta \log \pi_\theta(y \mid x)
\cdot
\frac{P_{\mathrm{expert}}(y \mid x)}{\pi_\theta(y \mid x)}
\right].
\end{aligned}
\label{eq:sft_importance}
\end{equation}

For deterministic expert demonstrations, the expert conditional distribution reduces to a Dirac delta:
\begin{equation}
P_{\mathrm{expert}}(y \mid x) = \mathbb{I}[y = y^*].
\end{equation}

Substituting this into Eq.~\eqref{eq:sft_importance} yields
\begin{equation}
\begin{aligned}
&\nabla_\theta \mathcal{L}_{\mathrm{SFT}}(\theta)\\
&=
-
\mathbb{E}_{(x,y) \sim Q}
\left[
\frac{\mathbb{I}[y = y^*]}{\pi_\theta(y \mid x)}
\nabla_\theta \log \pi_\theta(y \mid x)
\right].
\end{aligned}
\label{eq:sft_final_form_app}
\end{equation}

This recovers the equivalent on-policy formulation presented in the main text.

\subsection{Reinforcement Learning Interpretation}

Equation~\eqref{eq:sft_final_form_app} admits a direct reinforcement learning interpretation. In particular, it corresponds to an on-policy policy gradient with:
\begin{itemize}
    \item \textbf{Policy:} \(\pi_\theta(y \mid x)\);
    \item \textbf{Reward function:}
    \begin{equation}
    r(x,y) = \mathbb{I}[y = y^*],
    \end{equation}
    which provides a unit reward only when the sampled output exactly matches the expert demonstration;
    \item \textbf{Importance weight:}
    \begin{equation}
    w(x,y) = \frac{1}{\pi_\theta(y \mid x)},
    \end{equation}
    correcting for sampling from the model policy instead of the expert distribution.
\end{itemize}

Under this view, SFT can be regarded as a degenerate RL setting with an extremely sparse reward signal and high variance, where learning occurs only through trajectories that coincide exactly with expert demonstrations.

\subsection{Summary}

In summary, the derivation proceeds by (i) expressing the SFT gradient as an expectation over expert data, (ii) applying importance sampling to rewrite it under the model policy, and (iii) specializing the expert distribution to a deterministic form. This establishes a formal equivalence between SFT and on-policy reinforcement learning with a sparse indicator reward, providing a unified perspective on supervised and reinforcement-based post-training.

\section{Formulation of Group Fine-Tuning}
\label{sec:formulation}

In this appendix, we provide the explicit loss formulations and gradient expressions of Group Fine-Tuning (GFT), including both sequence-level and token-level forms. These formulations correspond to the gradient expression presented in Eq.~\eqref{eq:gft_gradient} in the main text.

\subsection{Sequence-Level Objective}
\label{subsec:seq_level_obj}

For each input query \(x\), we construct a response group
\(
\mathcal{G}_x = \{y_1, \ldots, y_K\}
\),
where each response \(y_k\) is assigned a scalar reward \(R(y_k)\) and a standardized
group advantage \(A(y_k)\) as defined in Eq.~\eqref{eq:advantage}. We define the sequence-level GFT loss as
\begin{equation}
\begin{aligned}
\mathcal{L}_{\mathrm{GFT}}^{\mathrm{seq}}(\theta)
=
-
\mathbb{E}_{x}
\Bigg[
\sum_{y_k \in \mathcal{G}_x}
A(y_k)\,
\mathcal{C}\!\left(\pi_\theta(y_k \mid x)\right)
\\[-1mm]
\quad\cdot
\log \pi_\theta(y_k \mid x)
\Bigg].
\end{aligned}
\label{eq:gft_loss_seq}
\end{equation}

where \(\mathcal{C}(\cdot)\) is the dynamic coefficient rectification function defined
in Eq.~\eqref{eq:rectification_factor}.

Taking the gradient of Eq.~\eqref{eq:gft_loss_seq} yields the sequence-level policy gradient:
\begin{equation}
\begin{split}
\nabla_\theta \mathcal{L}_{\mathrm{GFT}}^{\mathrm{seq}}
&= 
\mathbb{E}_{x}
\Bigg[
\sum_{y_k \in \mathcal{G}_x}
A(y_k)\,
\frac{\mathcal{C}\left(\pi_\theta(y_k \mid x)\right)}{\pi_\theta(y_k \mid x)}
\\
&\qquad\cdot 
\nabla_\theta \log \pi_\theta(y_k \mid x)
\Bigg].
\end{split}
\label{eq:gft_grad_seq}
\end{equation}

\subsection{Token-Level Decomposition}
\label{subsec:token_level_decomp}

Each response sequence \(y_k = (y_{k,1}, \ldots, y_{k,T_k})\) is generated autoregressively by the policy:
\begin{equation}
\pi_\theta(y_k \mid x)
=
\prod_{t=1}^{T_k}
\pi_\theta\!\left(y_{k,t} \mid y_{k,<t}, x\right).
\end{equation}

Accordingly, the sequence log-probability decomposes as
\begin{equation}
\log \pi_\theta(y_k \mid x)
=
\sum_{t=1}^{T_k}
\log \pi_\theta\!\left(y_{k,t} \mid y_{k,<t}, x\right).
\end{equation}

We use the shorthand
\begin{equation}
\pi_{k,t}
\triangleq
\pi_\theta\!\left(y_{k,t} \mid y_{k,<t}, x\right).
\label{eq:pi_kt_def}
\end{equation}
for the token-level prediction probability.

Substituting the above decomposition into Eq.~\eqref{eq:gft_loss_seq}, we obtain the token-level GFT loss:
\begin{equation}
\begin{split}
& \mathcal{L}_{\mathrm{GFT}}^{\mathrm{tok}}(\theta)
= % 第一行
\\
& - % 第二行 & 放在最前面，这样就会和第一行的开头对齐（顶格）
\mathbb{E}_{x}
\Bigg[
\sum_{y_k \in \mathcal{G}_x}
A(y_k)\,
\sum_{t=1}^{T_k}
\mathcal{C}(\pi_{k,t})\,
\log \pi_{k,t}
\Bigg].
\end{split}
\label{eq:gft_loss_tok}
\end{equation}

Taking the gradient yields the token-level policy gradient:
\begin{equation}
\begin{split}
& \nabla_\theta \mathcal{L}_{\mathrm{GFT}}^{\mathrm{tok}}
= % 第一行结束在等号
\\
& \mathbb{E}_{x} % & 放在最前面，确保与第一行开头对齐
\Bigg[
\sum_{y_k \in \mathcal{G}_x}
A(y_k)\,
\sum_{t=1}^{T_k}
\frac{\mathcal{C}(\pi_{k,t})}{\pi_{k,t}}\,
\nabla_\theta
\log \pi_{k,t}
\Bigg].
\end{split}
\label{eq:gft_grad_tok}
\end{equation}

\subsection{Relation to SFT and RL Objectives}
\label{subsec:relation_sft_rl}

When the response group degenerates to a single expert demonstration (\(|\mathcal{G}_x| = 1\)), the advantage is constant and Eq.~\eqref{eq:gft_grad_tok} reduces to the standard SFT gradient. Conversely, when the group consists of diverse sampled trajectories with non-trivial advantage values, GFT recovers an on-policy reinforcement learning update with group-normalized advantage weighting and bounded importance coefficients.

This formulation establishes GFT as a strict generalization of SFT and a stabilized, contrastive variant of policy-gradient-based post-training.

% \section{More Experiment}

% \subsection{Main Results on VLLM base model}\label{sec:main_results_VLLM}

% \begin{table*}[t]
%     \centering
%     \small % 稍微减小字号
%     \caption{Main results on Multimodal reasoning benchmarks. We compare GFT with various baselines across different model families. \textbf{Bold} indicates the best performance within each model group. Note that GFT and GRPO use only \textbf{10k} training data, while other methods use \textbf{100k}.}
%     \label{tab:main_results_mllm}
    
%     % 使用 resizebox 强制将表格宽度适应页面宽度 (\textwidth)
%     \resizebox{\textwidth}{!}{%
%         \setlength{\tabcolsep}{3pt} % 稍微减小列间距以适应更多列
%         \renewcommand{\arraystretch}{1.2} % 稍微增加行高
%         % 列数由 9 改为 8 (Model, Method, 6个Benchmark)
%         \begin{tabular}{llcccccc}
%             \toprule
%             \textbf{Model} & \textbf{Method} & \textbf{AMC23} & \textbf{College Math} & \textbf{Gaokao2023En} & \textbf{Math} & \textbf{Minerva Math} & \textbf{TabMWP} \\
%             \midrule
            
%             % 第一组: Qwen2.5-Math-1.5B
%             \multirow{7}{*}{\textbf{Qwen2.5-Math-1.5B}} 
%             & Base Model & 19.38 & - & - & 31.66 & 8.51 & - \\
%             & + SFT & 31.25 & 36.45 & 48.86 & 60.66 & 23.99 & 79.34 \\
%             & + GRPO & 44.84 & 35.58 & 51.80 & 65.97 & 21.17 & 76.94 \\
%             & + ASFT & 43.12 & 29.40 & 47.99 & 60.35 & 15.55 & 65.06 \\
%             & + PSFT & 31.56 & 33.77 & 47.66 & 59.51 & 19.13 & 71.61 \\
%             & + DFT & 36.40 & 38.76 & 52.75 & 64.35 & 23.75 & 82.08 \\
%             \cmidrule(l){2-8} % 横线范围调整为 2-8
%             & \textbf{+ GFT (Ours)} & \textbf{46.09} & \textbf{40.51} & \textbf{58.32} & \textbf{70.50} & \textbf{30.42} & \textbf{85.24} \\
%             \midrule
            
%             % 第二组: Qwen2.5-Math-7B
%             \multirow{7}{*}{\textbf{Qwen2.5-Math-7B}} 
%             & Base Model & - & - & - & - & - & - \\
%             & + SFT & 41.88 & 38.31 & 54.69 & 67.16 & 31.82 & 87.67 \\
%             & + GRPO & 56.09 & 39.81 & 63.47 & 74.52 & 37.22 & 92.36 \\
%             & + ASFT & 52.81 & 40.76 & 61.55 & 74.31 & 32.47 & 89.37 \\
%             & + PSFT & 41.56 & 38.05 & 56.74 & 67.30 & 34.86 & 83.55 \\
%             & + DFT & 51.09 & 39.31 & 57.46 & 70.42 & 35.31 & 86.94 \\
%             \cmidrule(l){2-8}
%             & \textbf{+ GFT (Ours)} & \textbf{59.47} & \textbf{42.06} & \textbf{65.81} & \textbf{77.31} & \textbf{39.86} & \textbf{93.81} \\
%             \midrule
            
%             % 第三组: DeepSeekMath-7B
%             \multirow{6}{*}{\textbf{DeepSeekMath-7B}} 
%             & Base Model & - & - & - & - & - & - \\
%             & + SFT & - & - & - & - & - & - \\
%             & + GRPO & - & - & - & - & - & - \\
%             & + ASFT & - & - & - & - & - & - \\
%             & + PSFT & - & - & - & - & - & - \\
%             & + DFT & - & - & - & - & - & - \\
%             \cmidrule(l){2-8}
%             & \textbf{+ GFT (Ours)} & - & - & - & - & - & - \\
%             \midrule
            
%             % 第四组: LLaMA-3.1-8B
%             \multirow{6}{*}{\textbf{LLaMA-3.1-8B}} 
%             & Base Model & - & - & - & - & - & - \\
%             & + SFT & - & - & - & - & - & - \\
%             & + GRPO & - & - & - & - & - & - \\
%             & + ASFT & - & - & - & - & - & - \\
%             & + PSFT & - & - & - & - & - & - \\
%             & + DFT & - & - & - & - & - & - \\
%             \cmidrule(l){2-8}
%             & \textbf{+ GFT (Ours)} & - & - & - & - & - & - \\
%             \bottomrule
%         \end{tabular}%
%     }
% \end{table*}

% \subsection{Learning Dynamic}\label{moredynamic}

% \subsection{Pass@k Accuracy and Diversity of GFT}\label{morepass}

% \begin{table}[t]
%     \centering
%     % 更新 Caption 中的 k 值为 8, 16 以匹配表格内容
%     \caption{Comparison of Pass@k ($k=8, 16$) performance between Distillation, GRPO and GFT. High Pass@k indicates that the model has learned a diverse and robust distribution of correct reasoning paths.}
%     \label{tab:passk_appendix}
%     \resizebox{\linewidth}{!}{%
%         \renewcommand{\arraystretch}{1.4}
%         \setlength{\tabcolsep}{5pt}
%         \begin{tabular}{c|l|ccccc}
%             \toprule
%             \textbf{Metric} & \textbf{Method} & \textbf{Math500} & \textbf{Minerva} & \textbf{Olympiad} & \textbf{AIME24} & \textbf{AMC23} \\
%             \midrule
%             % Pass@8 Section: multirow 改为 3 行
%             \multirow{3}{*}{Pass@8} 
%             & Distill & 00.00 & 00.00 & 00.00 & 00.00 & 00.00 \\
%             & GRPO & 00.00 & 00.00 & 00.00 & 00.00 & 00.00 \\
%             & \textbf{GFT (Ours)} & \textbf{00.00} & \textbf{00.00} & \textbf{00.00} & \textbf{00.00} & \textbf{00.00} \\
%             \midrule
%             % Pass@16 Section: multirow 改为 3 行
%             \multirow{3}{*}{Pass@16} 
%             & Distill & 00.00 & 00.00 & 00.00 & 00.00 & 00.00 \\
%             & GRPO & 00.00 & 00.00 & 00.00 & 00.00 & 00.00 \\
%             & \textbf{GFT (Ours)} & \textbf{00.00} & \textbf{00.00} & \textbf{00.00} & \textbf{00.00} & \textbf{00.00} \\
%             \bottomrule
%         \end{tabular}%
%     }
% \end{table}

% \subsection{Hyperparameter Analysis}\label{morehyper}

% \section{Training Settings}
% % \paragraph{Training Settings}
% Following DFT~\citep{wu2025generalization}, we utilize the NuminaMath CoT dataset~\citep{numina_math_datasets}, selected for its extensive diversity ranging from high school exercises to international olympiads. For GFT, we construct a hybrid response group of size $K=8$ per query, comprising 1 expert demonstration, 3 teacher distillations from Qwen2.5-Math-72B, and 4 self-generated samples. Similarly, the GRPO baseline is configured to generate 8 outputs per query. To align the total volume of training data across different paradigms, GFT and GRPO utilize a randomly sampled 10k subset (processing 8 trajectories per query), whereas single-trajectory baselines (e.g., SFT, DFT) are trained on a randomly sampled 100k subset.
% \section{Related work}

\section{Evaluation Settings}
\label{sec:eval_settings}
We conduct evaluations on a broad suite of 11 benchmarks: AMC23~\citep{amc2023dataset}, College Math~\citep{hendrycks2020measuring}, Gaokao~\citep{zhang2023evaluating}, Math~\citep{hendrycks2021measuring}, Minerva Math~\citep{lewkowycz2022solving}, TabMWP~\citep{lu2022dynamic}, OlympiadBench~\citep{he-etal-2024-olympiadbench}, Mmlu Stem~\citep{hendrycks2020measuring}, Sat Math~\citep{zhong2024agieval}, Mawps~\citep{koncel2016mawps}, and Svamp~\citep{patel2021nlp}. These benchmarks are carefully selected to cover a wide spectrum of difficulty levels and reasoning types, ensuring a holistic assessment of the model's capabilities. We report the average Pass@1 accuracy across 16 decoding runs (Pass@16 Average) with a sampling temperature of 0.5 and a maximum generation length of 4096 tokens.


% 第一段
\paragraph{Trade-off Between SFT and RL}
Post-training paradigms typically navigate a trade-off between Supervised Fine-Tuning (SFT) and Reinforcement Learning (RL). SFT is widely recognized for its efficiency in knowledge injection and ``cold-starting''~\citep{zhou2023lima, chung2024scaling}; however, it is prone to mechanical memorization and often fails to generalize to out-of-distribution scenarios~\citep{ouyang2022training, bai2022training, chu2025sft, swamy2025all, huan2025does}. Conversely, RL excels at discovering robust strategies and optimizing long-term objectives~\citep{christiano2017deep}, yet it is computationally expensive and struggles to acquire complex reasoning skills from scratch without sufficient guidance~\citep{schulman2017proximal, sheng2025hybridflow, mandlekar2021matters, chen2025empirical}.

% 第二段
\paragraph{The Synergy Dilemma in Hybrid Post-Training}
Standard hybrid approaches (e.g., SFT followed by RL) attempt to combine these complementary strengths but face a severe ``synergy dilemma''~\citep{ouyang2022training, rafailov2023direct}. Recent studies conclude that this conflict arises from the fundamental training dynamics: the overfitting induced by SFT creates a rigid policy that severely constrains the exploration space required for subsequent RL~\citep{chen2025sft}, while simultaneously leading to reasoning pattern mismatches that hinder effective policy alignment~\citep{chen2025synergy}. Although methods like interleaved updates~\citep{liu2025uft} or preference optimization~\citep{rafailov2023direct} offer partial solutions, they remain dependent on external feedback signals. In contrast, our work addresses this dilemma by transforming the rigid imitation objective into a \textbf{Group Advantage Learning} framework, which explicitly preserves solution diversity and the exploration manifold by optimizing contrastive advantages derived from hybrid response groups.

% 第三段
\paragraph{Single-Stage Hybrids: Mixing Imitation and Exploration}
Several recent studies have attempted to unify SFT and RL by balancing imitation and exploration through modified objectives~\citep{yuan2024self}. Single-stage hybrid methods, such as SRFT~\citep{fu2025srft} and UFT~\citep{liu2025uft}, employ dynamic weighting mechanisms, interleaved updates, or dense verification signals~\citep{wang2024math, yu2024ovm} to mix supervised signals with reinforcement objectives. Similarly, frameworks like HybridFlow~\citep{sheng2025hybridflow} explore flexible combinations of offline and online data to bridge the gap. While approaches like CHORD~\citep{zhu2025anchored} introduce anchor-based constraints to maintain stability, a common limitation across these methods is that they often treat SFT and RL as separate components to be linearly combined or alternated, rather than fusing them mathematically into a cohesive formulation derived from a unified training dynamic.

% 第四段
\paragraph{Gradient-Level Stabilization and Its New Trade-offs}
To address the instability inherent in post-training, other researchers have revisited the underlying gradient formulation. Theoretical analyses suggest a deeper equivalence between likelihood maximization and reinforcement learning~\citep{swamy2025all}, prompting new rectification strategies. For instance, \citet{wu2025generalization} propose Dynamic Fine-Tuning (DFT), which counteracts gradient explosion by reweighting the loss with the model's likelihood to cancel the inverse-probability term. However, this indiscriminate dampening creates a new dilemma: it suppresses the strong gradient signals required for injecting novel knowledge, potentially hindering adaptation to new domains. Alternatively, approaches like Proximal SFT~\citep{zhu2025proximal} and Anchored SFT~\citep{zhu2025anchored} introduce trust-region constraints to stabilize fine-tuning, yet such rigid regularizations may overly constrain the model's plasticity. In the realm of Reinforcement Learning, stability is traditionally enforced via KL-divergence penalties~\citep{ouyang2022training} or clipping mechanisms~\citep{schulman2017proximal}. More recently, group-based methods like GRPO~\citep{shao2024deepseekmath} have emerged to mitigate gradient variance by normalizing advantages within generated groups, effectively removing the reliance on unstable critic models, while system-level frameworks like HybridFlow~\citep{sheng2025hybridflow} attempt to stabilize training through flexible data scheduling.
